\section{Tillförlitlig leverans}
\label{sec:reliability}

Checksumman ger \textbf{integritet}: den som tar emot en frame vet att den är hel. Den ger ingen
\textbf{leveransgaranti}: sändaren vet fortfarande inte om något kom fram.

Skillnaden är hela det här avsnittet, och de fyra mekanismer som följer bygger på varandra i tur
och ordning.

\subsection{ACK och NACK}

Mottagaren svarar. Kom framen fram hel skickar den ett \code{Ack}; var den trasig på ett sätt som
ändå går att svara på skickar den ett \code{Nack}.

Reglerna för kvittensframen:

\begin{itemize}
  \item \code{DST} och \code{SRC} byter plats mot den kvitterade framen.
  \item \code{SEQ} är \textbf{samma} som i den kvitterade framen --- det är så sändaren vet vad
    som kvitteras.
  \item \code{LEN} är noll.
\end{itemize}

Att \code{SEQ} följer med är inte en detalj. Utan det vet en sändare som har två frames i luften
inte vilken av dem som kom fram.

\begin{aside}
\note Ett \code{Nack} går bara att skicka när framen är trasig på ett sätt som lämnar avsändare
och sekvensnummer läsbara. Är checksumman fel kan man inte lita på \emph{något} fält --- inklusive
\code{SRC} --- så i praktiken är utebliven kvittens det vanliga sättet ett fel visar sig.
\end{aside}

\subsection{Timeout}

Och därmed: vad gör sändaren när ingenting kommer tillbaka?

Den kan inte vänta för alltid, och den kan inte veta skillnaden mellan en förlorad frame, en
förlorad kvittens och en långsam mottagare. Den enda utvägen är att bestämma sig efter en viss
tid: en \textbf{timeout}.

Att välja timeout-värde är en avvägning utan rätt svar:

\begin{narrowtable}{@{}ll@{}}
\thead{För kort} & \thead{För lång} \\ \midrule
Onödiga omsändningar & Långsam felupptäckt \\
Fler dubbletter & Bussen står still och väntar \\
Bussen belastas mer när den redan är belastad & Systemet känns hängt \\
\end{narrowtable}

\subsection{Omsändning}

Vid timeout skickas framen om, med \textbf{samma} \code{SEQ}. Det är avgörande: ett nytt
sekvensnummer skulle göra omsändningen till ett nytt meddelande, och mottagaren skulle inte kunna
känna igen den som en upprepning.

Antalet försök måste vara begränsat. En sändare som försöker i all evighet mot en nod som är
avstängd gör ingenting annat.

\subsection{Dubbletter}

Och här stänger cirkeln. Omsändningen är precis det som skapar dubbletter.

Det farliga fallet är inte att framen tappas --- då hände ingenting --- utan att \textbf{kvittensen}
tappas:

\begin{console}
A -> B   StatusRequest, SEQ 0x0007      B tar emot och behandlar den
B -> A   Ack, SEQ 0x0007                kvittensen tappas
A        timeout
A -> B   StatusRequest, SEQ 0x0007      omsändning
B                                       ...behandlar den en andra gång?
\end{console}

Ur B:s synvinkel är det två identiska frames. Behandlar den båda har kommandot utförts två gånger,
och är kommandot "öka börvärdet med ett" blir resultatet fel.

\textbf{Lösningen är att komma ihåg.} Varje nod sparar senast behandlade \code{SEQ} per avsändare.
Kommer samma \code{SEQ} igen:

\begin{itemize}
  \item Kör \textbf{inte} applikationslogiken igen.
  \item Skicka \textbf{ändå} kvittensen igen.
\end{itemize}

Den andra punkten är den som glöms. Sändaren väntar fortfarande på sin kvittens, och tystnad
leder bara till ännu en omsändning.

\begin{aside}
\note Fyra mekanismer, i en kedja där var och en löser problemet den föregående skapade:
checksumma upptäcker fel, ACK bekräftar leverans, timeout upptäcker utebliven ACK, omsändning
åtgärdar den --- och dubblettdetektering städar upp efter omsändningen. Ingen av dem går att
utelämna, och det är därför kedjan är värd att kunna utantill.
\end{aside}

\subsection{Och sedan gör hårdvaran det åt oss}

Allt i det här avsnittet är kod du skriver, testar och underhåller. Från \kapref{ch:can} ligger
det mesta av det i kisel: CAN gör ramning, checksumma och kollisionshantering i hårdvara, och
kvitterar dessutom i en enda bit inuti framen.

Det gör inte kapitlet mindre viktigt. Det gör det till den bakgrund mot vilken CAN går att
förstå --- och till en del av vad duggan examinerar.
